% GENERATED FILE. Edit canonical lessons, then run npm run generate:books.
% canonical-source-hash: fnv1a64:7dde870aa4a56f28
% canonical-lessons: ES-C560-congregarse, ES-C560-juntarse, ES-C560-convocar, ES-C560-tertulia, ES-C560-reencuentro

\chapter{El reencuentro}
\label{ch:es-el-reencuentro}
\hlchaptermodality{560}

\begin{chapteropening}
\textbf{By the end of this chapter:} \emph{I can say that people gathered, got together, met to talk, were called to a meeting and saw each other again.}

Complete the last lesson of chapter 560: I can say that people gathered, got together, met to talk, were called to a meeting and saw each other again.
\end{chapteropening}

\section[congregarse]{congregarse — to gather, to assemble}

\label{lesson:ES-C560-congregarse}



\begin{quote}
Before the new one: say the Spanish for the float, then the Spanish for the crowd.
\end{quote}

\subsection*{You'll want to know: congregarse}
\textbf{congregarse} — \textquotedblleft{}to gather, to assemble\textquotedblright{}.

\begin{quote}
Cientos de personas se congregaron en la plaza.
\end{quote}

\begin{quote}
\emph{Hundreds of people gathered in the square.}
\end{quote}

It takes \emph{gu} before \emph{e}: \emph{se congreguen}. Common in the news.

\begin{cousinweb}
\emph{congregarse} is a learned borrowing from Latin \emph{congregāre} 'to gather into a flock, to assemble' (Lewis \& Short), built on \emph{grex}, \emph{gregis} 'a flock'.
\end{cousinweb}

\begin{grammarlens}[title={What you've built}]
One more word for telling, explaining and describing at B1.
\end{grammarlens}

\subsection*{Your turn}
\begin{itemize}
\raggedright
  \item \emph{Say it:} \emph{congregarse}
  \item \emph{Say it:} \emph{congregarse}, once more
  \item \emph{Say it:} a sentence of your own with \emph{congregarse}
  \item \emph{Recall:} say the Spanish for the float, then the Spanish for the crowd, then say \emph{congregarse} again
\end{itemize}

\begin{tcolorbox}[breakable,colback=teal!4,colframe=teal!35!black,title={Before you move on}]
What does congregarse mean? (\textquotedblleft{}To gather, to assemble\textquotedblright{}.) Say it once more.
\end{tcolorbox}
\section[juntarse]{juntarse — to get together; to join}

\label{lesson:ES-C560-juntarse}



\begin{quote}
Before the new one: say the Spanish for the crowd, then the Spanish for to gather.
\end{quote}

\subsection*{You'll want to know: juntarse}
\textbf{juntarse} — \textquotedblleft{}to get together; to join\textquotedblright{}.

\begin{quote}
Los amigos se juntan los viernes.
\end{quote}

\begin{quote}
\emph{The friends get together on Fridays.}
\end{quote}

Common in Latin America for meeting friends. \emph{Juntarse con alguien} — to hang out with someone.

\begin{cousinweb}
\emph{juntarse} is built on \emph{junto} 'together', which comes from Latin \emph{iūnctus} 'joined', the participle of \emph{iungere} 'to join, to yoke' (Lewis \& Short).
\end{cousinweb}

\begin{grammarlens}[title={What you've built}]
One more word for telling, explaining and describing at B1.
\end{grammarlens}

\subsection*{Your turn}
\begin{itemize}
\raggedright
  \item \emph{Say it:} \emph{juntarse}
  \item \emph{Say it:} \emph{juntarse}, once more
  \item \emph{Say it:} a sentence of your own with \emph{juntarse}
  \item \emph{Recall:} say the Spanish for the crowd, then the Spanish for to gather, then say \emph{juntarse} again
\end{itemize}

\begin{tcolorbox}[breakable,colback=teal!4,colframe=teal!35!black,title={Before you move on}]
What does juntarse mean? (\textquotedblleft{}To get together; to join\textquotedblright{}.) Say it once more.
\end{tcolorbox}
\section[convocar]{convocar — to call (a meeting, a strike)}

\label{lesson:ES-C560-convocar}



\begin{quote}
Before the new one: say the Spanish for to gather, then the Spanish for to get together.
\end{quote}

\subsection*{You'll want to know: convocar}
\textbf{convocar} — \textquotedblleft{}to call (a meeting, a strike)\textquotedblright{}.

\begin{quote}
Los trabajadores convocaron una huelga.
\end{quote}

\begin{quote}
\emph{The workers called a strike.}
\end{quote}

It takes \emph{qu} before \emph{e}: \emph{convoqué}.

\begin{cousinweb}
\emph{convocar} is a learned borrowing from Latin \emph{convocāre} 'to call together' (Lewis \& Short), built on \emph{con-} and \emph{vocāre} 'to call'.
\end{cousinweb}

\begin{grammarlens}[title={What you've built}]
One more word for telling, explaining and describing at B1.
\end{grammarlens}

\subsection*{Your turn}
\begin{itemize}
\raggedright
  \item \emph{Say it:} \emph{convocar}
  \item \emph{Say it:} \emph{convocar}, once more
  \item \emph{Say it:} a sentence of your own with \emph{convocar}
  \item \emph{Recall:} say the Spanish for to gather, then the Spanish for to get together, then say \emph{convocar} again
\end{itemize}

\begin{tcolorbox}[breakable,colback=teal!4,colframe=teal!35!black,title={Before you move on}]
What does convocar mean? (\textquotedblleft{}To call (a meeting, a strike)\textquotedblright{}.) Say it once more.
\end{tcolorbox}
\section[la tertulia]{la tertulia — the get-together (for conversation)}

\label{lesson:ES-C560-tertulia}



\begin{quote}
Before the new one: say the Spanish for to get together, then the Spanish for to call.
\end{quote}

\subsection*{You'll want to know: la tertulia}
\textbf{la tertulia} — \textquotedblleft{}the get-together (for conversation)\textquotedblright{}.

\begin{quote}
Los viernes hay tertulia en el café.
\end{quote}

\begin{quote}
\emph{On Fridays there is a get-together at the café.}
\end{quote}

A regular group that meets to talk; on radio and TV, \emph{una tertulia} is a panel discussion.

\begin{cousinweb}
The origin of \emph{tertulia} is not settled. A common account links it to the Latin writer Tertullian, read and argued over in old gatherings, but this is not proved.
\end{cousinweb}

\begin{grammarlens}[title={What you've built}]
One more word for telling, explaining and describing at B1.
\end{grammarlens}

\subsection*{Your turn}
\begin{itemize}
\raggedright
  \item \emph{Say it:} \emph{la tertulia}
  \item \emph{Say it:} \emph{la tertulia}, once more
  \item \emph{Say it:} a sentence of your own with \emph{la tertulia}
  \item \emph{Recall:} say the Spanish for to get together, then the Spanish for to call, then say \emph{la tertulia} again
\end{itemize}

\begin{tcolorbox}[breakable,colback=teal!4,colframe=teal!35!black,title={Before you move on}]
What does la tertulia mean? (\textquotedblleft{}The get-together (for conversation)\textquotedblright{}.) Say it once more.
\end{tcolorbox}
\section[el reencuentro]{el reencuentro — the reunion}

\label{lesson:ES-C560-reencuentro}



\begin{quote}
Before the new one: say the Spanish for to call, then the Spanish for the get-together.
\end{quote}

\subsection*{You'll want to know: el reencuentro}
\textbf{el reencuentro} — \textquotedblleft{}the reunion\textquotedblright{}.

\begin{quote}
El reencuentro de los hermanos fue muy bonito.
\end{quote}

\begin{quote}
\emph{The brothers' reunion was lovely.}
\end{quote}

\begin{cousinweb}
\emph{reencuentro} is built on \emph{encuentro} with \emph{re-} 'again'. \emph{Encuentro} comes from \emph{encontrar}, which EtymDB traces to Latin \emph{in contrā} 'facing'.
\end{cousinweb}

\begin{grammarlens}[title={What you've built}]
One more word for telling, explaining and describing at B1.
\end{grammarlens}

\subsection*{Your turn}
\begin{itemize}
\raggedright
  \item \emph{Say it:} \emph{el reencuentro}
  \item \emph{Say it:} \emph{el reencuentro}, once more
  \item \emph{Say it:} a sentence of your own with \emph{el reencuentro}
  \item \emph{Recall:} say the Spanish for to call, then the Spanish for the get-together, then say \emph{el reencuentro} again
\end{itemize}

\begin{tcolorbox}[breakable,colback=teal!4,colframe=teal!35!black,title={Before you move on}]
What does el reencuentro mean? (\textquotedblleft{}The reunion\textquotedblright{}.) Say it once more.
\end{tcolorbox}
